\begin{theorem}[independent certification of every box]\label{push:src:E.4.1}


For all 19701 rational rectangles in the supplied certificate, and for
each $\tau$ in the rectangle, both zeros of $\wp_\tau$ satisfy
$\mathcal N^2<2$. The new largest certified upper bound is less than
\[
 1.999877571767410525.
\]
The 19701 rectangles, 36 local rectangles and 174 excluded rectangles
have disjoint interiors and cover exactly
$[0,1/2]\times[17/20,2]$. Every local rectangle is contained in
$|\tau-i|<3/20000$, and every excluded rectangle satisfies
$|\tau|<1$ throughout.
\end{theorem}



\begin{proof}
\begin{enumerate}[label=\arabic*.,leftmargin=*]
\item The programs \path{code/I/replay_arb/S6_02_PUSH/code/e4_arb_certificate.py} and
\path{code/I/replay_arb/S6_02_PUSH/code/e4_cover_and_local.py} use respectively 192-bit and 256-bit
Arb arithmetic through python-flint 0.9.0. The arithmetic implementation
is independent of mpmath. We import neither mpmath nor any routine
from the previous certificate. The only numerical data taken from that
certificate are its rational rectangle endpoints. We recompute roots,
preconditioners, radii, contraction bounds, function values and tails.
Arb's midpoint-radius representation and its outward error bounds are
described by \citet[Section~2, pp.~2--4, and Section~3, pp.~4--6]{Johansson2016}; the precise
constructor and endpoint semantics used here are documented in
\citet{PythonFlint2026}. The final comparison is made with an upper
endpoint, before conversion to a decimal string.

\item For a rectangle $Q$, let $z_0$ and $A\ne0$ be the exact dyadic complex
numbers selected by the midpoint Newton iteration. That iteration is
only a choice of constants. For $Z=z_0+[-r,r]+i[-r,r]$ we independently
enclose $\wp_\tau(z_0)$ and $\wp'_\tau(Z)$ for $\tau\in Q$.
We accept only if
\[
 \begin{split}
 &\left|-A\wp_\tau(z_0)
 +(1-A\wp'_\tau(Z))(Z-z_0)\right|_\infty<r,\\
 &|\Re(1-A\wp'_\tau(Z))|
       +|\Im(1-A\wp'_\tau(Z))|<1.
 \end{split}
\]
Here absolute values on enclosures mean their outward upper bounds.
The derivative enclosure and integration along a line segment show
that $z\mapsto z-A\wp_\tau(z)$ is a strict contraction of $Z$ into
itself, uniformly for $\tau\in Q$. Thus it has a genuine zero in $Z$.
The finite Fourier enclosures exclude the poles on $Z$, and the strict
Lipschitz inequality forces $A\ne0$. The two zeros of $\wp$ are
negatives modulo the period lattice, counted with multiplicity. The
exact half-period transformation, obtained by replacing the summation
index in the theta series, is
\[
 \vartheta_3\bigl(z-(1-\tau)/2,\tau\bigr)
   =i e^{-\pi i\tau/4-\pi i z}\vartheta_1(z,\tau).
\]
Writing $v=\Im z$ and $y=\Im\tau$, its modulus and Gaussian factor
cancel to give
\[
 \mathcal N(z,\tau)=
 \left|\frac{\vartheta_1(z,\tau)}{\eta_D(\tau)}\right|
 e^{-\pi v^2/y}.
\]
Now $\vartheta_1(-z)=-\vartheta_1(z)$,
$\vartheta_1(z+1)=-\vartheta_1(z)$ and
$\vartheta_1(z+\tau)=-e^{-\pi i\tau-2\pi i z}\vartheta_1(z)$.
In the last transformation the modulus contributes
$e^{\pi y+2\pi v}$, while the Gaussian contributes its reciprocal.
Thus the displayed norm is even and lattice-periodic, proving directly
that the one-root enclosure controls both zeros. No continuity or
success of Newton iteration is used in that conclusion.

\item We use the normally convergent cosine expansion
\[
 \wp_\tau(z)=\frac{\pi^2}{\sin^2\pi z}-\frac{\pi^2}{3}
 +8\pi^2\sum_{n\geq1}\frac{nq^n}{1-q^n}
                      (1-\cos 2\pi nz),\qquad q=e^{2\pi i\tau},
\]
and differentiate it in $z$ for $\wp'$. This identity follows by twice
differentiating the normally convergent Jacobi product and inserting
$E_2=1-24\sum_{n\geq1}nq^n/(1-q^n)$; this also verifies the
Fourier formula used in the period calculation above. The new program sums
32 terms directly in cosine form. Put $r_q\geq|q|$,
$H_z\geq|\Im z|$, and $R_q=r_q e^{2\pi H_z}<1$.
For $j=1,2$, the inequality $(N+k)^j\leq(N+1)^jk^j$ gives
\[
 \begin{aligned}
 U_1(a)&=\frac{(N+1)a^{N+1}}{(1-a)^2},\\
 U_2(a)&=\frac{(N+1)^2a^{N+1}(1+a)}{(1-a)^3},
 \qquad
 \sum_{n>N}n^j a^n\leq U_j(a).
 \end{aligned}
\]
We therefore add complex rectangles enclosing discs of radii
\[
 \frac{8\pi^2}{1-r_q}\bigl(U_1(r_q)+U_1(R_q)\bigr),
 \qquad \frac{16\pi^3}{1-r_q}U_2(R_q)
\]
to the two partial sums. These estimates follow from
$|\cos z|,|\sin z|\leq e^{|\Im z|}$ and
$|1-q^n|\geq1-r_q$. They differ from the moment formulas used in
the old implementation.

\item For $\vartheta_3$ we sum
$1+2\sum_{n=1}^{12}e^{\pi i n^2\tau}\cos(2\pi nh)$.
With $T_*\leq\Im\tau$ and $H_*\geq|\Im h|$, the remaining
absolute value is at most
\[
 \frac{2e^{-169\pi T_*+26\pi H_*}}
 {1-e^{-27\pi T_*+2\pi H_*}}.
\]
The ratio in this expression is tested to be less than one. For eta
we independently sum its logarithm,
\[
 \log\eta_D(\tau)=\frac{\pi i\tau}{12}
          +\sum_{n=1}^{24}\log(1-q^n)+\mathcal E,
 \qquad
 |\mathcal E|\leq
 \frac{r_q^{25}}{(1-r_q)(1-r_q^{25})}.
\]
The logarithms are holomorphic on $|q^n|<1$. The tail follows by
summing $|\log(1-z)|\leq|z|/(1-|z|)$.
We enclose $\mathcal N^2$ as
\[
 ((\Re\vartheta_3)^2+(\Im\vartheta_3)^2)
 \exp\left(-2\Re\log\eta_D-
                   \frac{2\pi(\Im h)^2}{\Im\tau}\right),
\]
and test its upper endpoint against the exact integer 2.

\item All 19701 tests succeeded. The full output records each new enclosure;
the largest is enclosed by
\[
 1.999877571767410524070536478244687423578792301873142913
       \ \pm 9.26\cdot10^{-56}.
\]
This is the enclosure of a certified upper bound, not an estimate of
the actual maximum. Its box has input index 17866, with zero-based
indexing. The file digest of the input is recorded in the output.

\item The cover check is exact rational arithmetic. At all 438 distinct
vertical edges we maintain the active closed vertical intervals. In
every open vertical slab they form an ordered partition of
$[17/20,2]$, each with multiplicity one. The last active set is empty;
all leaves lie in the root rectangle, and the sum of their exact areas
is $23/40$. The closed rectangles also cover their common edges.
The omitted and excluded conditions are checked separately using
rational endpoint comparisons. Finally,
$2(1/10000)^2<(3/20000)^2$ puts the entire omitted square inside the
analytic disc used below. This verifies both the inequalities and
the coverage, rather than just replaying the old list of successful
flags.
\end{enumerate}
\end{proof}

